\chapter{آمادگی قبل از سفر}
\label{ch:preparation}

\section{مدارک و اسناد}

\begin{checklistbox}[title={مدارکی که باید همراه داشته باشید}]
  \begin{checklist}
    \item پاسپورت + کپی تمام صفحات
    \item ویزا (برچسب یا \lr{eVisa} پرینت‌شده)
    \item نامه پذیرش و \lr{LOA/CAS/I-20}
    \item مدارک مالی (نسخه اصلی یا کپی تأییدشده)
    \item بیمه‌نامه سفر و درمان
    \item ریزنمرات و ترجمه (نسخه اصلی)
    \item عکس‌های پرسنلی اضافه
  \end{checklist}
\end{checklistbox}

\section{مسکن}

\begin{itemize}
  \item \textbf{خوابگاه}: معمولاً برای ترم اول راحت‌تر است؛ زود درخواست دهید
  \item \textbf{آپارتمان مشترک}: ارزان‌تر؛ از گروه‌های دانشجویی بپرسید
  \item \textbf{موقت}: \lr{Airbnb} یا هتل برای هفته اول تا پیدا کردن خانه
\end{itemize}

\begin{tipbox}
  قبل از پرداخت اجاره از راه دور، از اعتبار آگهی مطمئن شوید.
  در صورت امکان از دانشجویان همان شهر نظر بگیرید.
\end{tipbox}

\section{امور مالی}

\begin{enumerate}
  \item مقدار محدودی پول نقد برای روزهای اول
  \item کارت بانکی بین‌المللی (\lr{Visa/Mastercard})
  \item افتتاح حساب محلی (معمولاً بعد از ورود)
  \item آشنایی با نرخ تبدیل و کارمزد حواله
\end{enumerate}

\section{بیمه}

بیمه درمانی در بسیاری از کشورها اجباری است.
گاهی دانشگاه بیمه ارائه می‌دهد؛ گاهی باید خودتان بخرید.
مدت پوشش را طوری انتخاب کنید که از روز ورود شروع شود.

\section{بسته‌بندی}

\begin{table}[htbp]
  \centering
  \caption{چک‌لیست بسته‌بندی}
  \label{tab:packing}
  \begin{tabular}{@{}p{0.45\textwidth}p{0.45\textwidth}@{}}
    \toprule
    \textbf{حتماً ببرید} & \textbf{بهتر است نبرید} \\
    \midrule
    مدارک و کپی‌ها & وسایل برقی ۲۲۰ ولت (بدون مبدل) \\
    داروهای شخصی + نسخه & مبلغ زیاد نقد \\
    لباس مناسب آب‌وهوا & کتاب درسی (اغلب دیجیتال است) \\
    چند وعده غذای فوری & وسایل با وزن زیاد \\
    \bottomrule
  \end{tabular}
\end{table}
